\documentclass{beamer}
\usetheme{CambridgeUS}
\usecolortheme{seahorse}
\setbeamertemplate{navigation symbols}{}

\usepackage[T2A]{fontenc}
\usepackage[russian]{babel}

\usepackage{mathpartir,stmaryrd}
\DeclareMathOperator{\elim}{elim}
\DeclareMathOperator{\type}{type}
\newcommand{\letin}[3]{\operatorname{let}\,#1=#2\,\operatorname{in}\,#3}

\title[DTSPL]{Системный язык программирования с зависимыми типами}
\author[Соколов П.П.]{Соколов П.П.\texorpdfstring{\\[0.3em]}{}
  {\small Научный руководитель Кузнецов С.Л.}}
\institute{МФТИ}
\date{21 июня 2024}

\begin{document}

\frame{\titlepage}

\begin{frame}{Мотивация}
  \begin{block}{Требования к системному программированию}
    \begin{minipage}[t]{0.48\linewidth}
    \begin{itemize}
      \item корректность;
      \item быстродействие;
      \item экономность по памяти;
    \end{itemize} 
    \end{minipage}
    \begin{minipage}[t]{0.48\linewidth}
    \begin{itemize}
      \item низкая энергозатратность;
      \item безопасность времени исполнения.
    \end{itemize}
    \end{minipage}
  \end{block}
  \begin{block}{Существующие СЯП}
  \begin{center}
  C/C++\qquad Zig\qquad Rust\qquad ATS\qquad \ldots
  \end{center}
  \end{block}
\end{frame}

\begin{frame}{Система типов Мартина-Лёфа}
  \begin{block}{Логические суждения}
  \begin{center}
  \begin{tabular}{rl}
    $\Gamma \vdash \tau \type$ & Формулировка типов \\
    $\Gamma \vdash t : \tau$ & Введение и устранение термов \\
    $\Gamma \vdash \sigma \equiv \tau$ & Равенство типов \\
    $\Gamma \vdash t \equiv u : t$ & Равенство термов \\
    $\Gamma \vdash$ & Корректность контекста
  \end{tabular}
  \end{center}
  \end{block}
  \begin{block}{Важнейшие правила вывода}
  \begin{mathpar}
    \inferrule
      {\Gamma,x:\sigma\vdash t:\tau}
      {\Gamma\vdash\lambda(x:\sigma).t:\Pi_{x:\sigma}\tau}\and
    \inferrule
      {\Gamma\vdash f:\Pi_{x:\sigma}\tau \\ \Gamma\vdash t:\sigma}
      {\Gamma\vdash f t:\tau[t/x]}\and
    \inferrule
      {\Gamma\vdash t:\mathbb{N}}
      {\Gamma\vdash S t:\mathbb{N}}\and
    \inferrule
      {\vdash t:\mathbb{N}
      \\x:\mathbb{N}\vdash\sigma\type
      \\\vdash t_0:\sigma[0/x]
      \\x:\mathbb{N},h:\sigma\vdash t_S:\sigma[S x/x]}
      {\vdash\elim^\mathbb{N}_{x.\sigma}t(0\Mapsto t_0;S x,h\Mapsto t_S)
        :\sigma[t/x]}
  \end{mathpar}
  \end{block}
\end{frame}

\begin{frame}{Линейная система типов с зависимыми типами}
  \begin{block}{Логические суждения}
  \begin{center}
  \begin{tabular}{rl}
    $\Delta; \Xi \vdash$ & Корректность контекста \\
    $\Delta; \cdot \vdash \tau \type$ & Формулировка типов \\
    $\Delta; \Xi \vdash t: \tau$ & Введение и устранение термов \\
    \ldots & Равенства контекстов, типов, термов
  \end{tabular}
  \end{center}
  \end{block}
  \begin{block}{Характерные правила вывода}
  \begin{mathpar}
    \inferrule
      {\vdash \Delta;\cdot\\(x:\tau)\in\Delta}
      {\Delta;\cdot\vdash x:\tau}\and
    \inferrule
      {\Delta;\cdot\vdash t:\sigma\\\Delta;\Xi\vdash u:\tau[t / x]}
      {\Delta;\Xi\vdash\;!t\otimes u:\Sigma_{!x:!\sigma}\tau}\and
    \inferrule
      {\Delta;\cdot\vdash\rho\type\\
      \Delta;\Xi\vdash t:\Sigma_{!x:!\sigma}\tau\\
      \Delta,x:\sigma;\Xi',y:\tau\vdash r:\rho}
      {\Delta;\Xi,\Xi'\vdash\letin{!x\otimes y}{t}{r}:\rho}
  \end{mathpar}
  \end{block}
\end{frame}

\begin{frame}{Система типов с часами}
\end{frame}

\begin{frame}{Итоговое исчисление}
\end{frame}

\begin{frame}{Метатеоремы}
\end{frame}

\begin{frame}{Дальнейшие расширения}
\end{frame}

\begin{frame}{Заключение}
\end{frame}

\end{document}
